% Transcription — fonds Alexandre Grothendieck, shelfmark 31, batch 5 (pages 81-100).
% Established from the facsimile archives/batches/31/batch-05.pdf (a mirror of
% https://grothendieck.umontpellier.fr/31.pdf), per the skill
% transcribe-grothendieck.
%
% Pass: Opus 5 (claude-opus-5), 2026-09-11 — first pass, unchecked against
% the pages by a human.
%
% The batch holds the end of the cd_ell run of batch 4, a page of conjectures,
% three further leaves of the corrected typescript begun on page 71, a short
% run on the Lefschetz theorem, and a long run opened by a cover sheet in his
% hand.
%
%   Pages 81-82 — ink: a Question on cd_ell(X) for X noetherian affine with the
%     counterexample that closes it, then « Dim. coh. d'un schéma affine sur Z »
%     — the Leray spectral sequence for X → Spec(Z), the two bounds on
%     R^q f_*(F) at the generic and at a closed point, and the boxed
%     cd_ell(X) ≤ dim X + 1.
%   Page 84 — ink, heavily worked over: « Conjectures sur la cohomologie des
%     schémas affines », 1°), 2°), 1 bis), 4°), with the invariant ν(k) defined
%     by ell^{ν(k)} = [k : k^ell].
%   Pages 83, 85, 87 — leaves of the same typescript as page 71 (batch 4),
%     scanned upside down and cancelled by long diagonals, the group letters and
%     several subscripts inked in his hand where the machine left gaps:
%     Théorème 3.1 (Hochschild cohomology of a group of multiplicative type),
%     Théorème 3.2, Corollaire 4.6 (schematic density), Corollaire 4.7,
%     Corollaire 4.9 and Remarque 4.9. Insertions in \add{}, deletions in
%     \struck{}.
%   Page 86 — a short N.B. in his hand on the verso of a typescript leaf.
%   Pages 89-90 — ink, fair: the Gysin sequence of a regular hyperplane section,
%     the Lefschetz theorem, H^{2n}(X, μ_N^{⊗n}) ≃ Z/NZ, and Corollaire 2.
%   Page 91 — an otherwise blank sheet inscribed « Théorème de Spécialisation :
%     Corollaire » in his hand; it opens the run of pages 92-98 and takes no
%     \page{}.
%   Pages 92-98 — ink, a fast draft: Lemme 1 (étale extensions of a local ring
%     of dimension ≥ 3), Lemme 2 (a smooth morphism with fibres of dimension 1
%     and a section), Lemme 3 (a regular local ring of dimension 2 and the
%     Abhyankar computation B = A[T]/(T^n − x)), the Théorème that a smooth
%     morphism is locally 1-aspherical for finite groups of order prime to the
%     residue characteristics, its Applications 1) 2) 3), and a Corollaire
%     listing three conditions on R^i h_*, R^i g_*.
%   Page 100 — an N.B. in his hand on the verso of an unrelated typed leaf,
%     on the spectral sequence R^p g_*(R^q i_*(G)) and R^i h_*(G).
%
% Page 88 (blank) is absent. Page 99 carries a leaf of an unrelated typescript
% (Bourbaki, Algèbre commutative, chap. VII, valuations essentielles d'un anneau
% de Krull) with no note of his, and is absent.
%
% Pages 92-98 are written at speed and their connective prose is largely lost;
% the statements and the formulas carry, and the apparatus says so rather than
% filling them in.
\documentclass[11pt,a4paper]{article}
% Le préambule commun à toutes les transcriptions du fonds Grothendieck.
%
% Il tient en une idée : **l'incertitude se compose, elle ne se résout pas.**
% Une transcription qui lisse un mot douteux a détruit la seule chose pour
% laquelle on la faisait — un lecteur doit pouvoir savoir, à chaque endroit,
% ce qui était sur le feuillet et ce qui a été ajouté. Les six macros qui
% suivent sont donc l'appareil critique, et non de la décoration.
%
% Les mêmes commandes sont reconnues par `scripts/render.mjs`, qui produit la
% vue de lecture du site. Une macro ajoutée ici doit l'être aussi là-bas,
% faute de quoi le rendu échouera bruyamment — ce qui est voulu : un rendu
% silencieusement faux est pire qu'un rendu absent.

\ProvidesPackage{grothendieck}[2026/08/08 Transcriptions du fonds Grothendieck]

\usepackage{amsmath,amssymb,amsthm}
\usepackage{mathrsfs}
\usepackage{stmaryrd}
% Les diagrammes commutatifs sont partout dans ce fonds : c'est la langue
% même de la géométrie algébrique de Grothendieck, et une transcription qui
% les rendrait en prose ne transcrirait rien.
\usepackage{tikz-cd}
% Français par défaut — c'est la langue des feuillets — mais l'anglais est
% chargé aussi : la traduction et le résumé sont des documents anglais, et la
% typographie française (espaces avant les deux-points) y serait une faute.
% Ces documents basculent par \selectlanguage{english} en tête de corps.
\usepackage[english,french]{babel}
\usepackage{fontspec}
\usepackage{geometry}
\geometry{a4paper,margin=25mm,marginparwidth=28mm}
\usepackage{marginnote}
\usepackage{xcolor}
% ulem, not soul. `soul`'s \st and \ul are built on a character-by-character
% scan that XeLaTeX's font machinery breaks: the document fails with an
% undefined control sequence rather than degrading. ulem does the same two jobs
% and is Unicode-engine safe. `normalem` stops it hijacking \emph, which the
% transcriptions use for Grothendieck's own emphasis.
\usepackage[normalem]{ulem}
\usepackage{hyperref}
\hypersetup{colorlinks=true,linkcolor=black,urlcolor=black,pdfborder={0 0 0}}

% --- Métadonnées du lot -----------------------------------------------------
% Reprises telles quelles par le script de rendu, qui les affiche en tête de
% la vue de lecture. Elles fixent ce que le fichier prétend couvrir : sans
% elles, une transcription flotte sans point d'attache dans un fonds de seize
% mille feuillets.

\newcommand{\folder}[1]{\gdef\grothFolder{#1}}
\newcommand{\batch}[1]{\gdef\grothBatch{#1}}
\newcommand{\leaves}[2]{\gdef\grothFirst{#1}\gdef\grothLast{#2}}
\newcommand{\foldertitle}[1]{\gdef\grothFoldertitle{#1}}
\gdef\grothFolder{?}\gdef\grothBatch{?}\gdef\grothFirst{?}\gdef\grothLast{?}\gdef\grothFoldertitle{}

% --- Le résumé --------------------------------------------------------------
%
% Il ouvre la lecture modernisée, sous le titre « Résumé », et il est composé
% en petit. Ce n'est pas un ornement : c'est le seul passage du document qui
% ne soit pas des mathématiques, et un lecteur doit pouvoir le reconnaître
% avant de l'avoir lu — puis le sauter s'il connaît déjà le sujet, ou s'y
% arrêter s'il ne le connaît pas. Le filet qui le suit marque la reprise du
% corps.

\newenvironment{resume}
  {\small\setlength{\parskip}{0.45em}}
  {\par\medskip\hrule\medskip}

% --- Mots-clés ---------------------------------------------------------------
%
% En anglais, en clôture du résumé de la lecture modernisée : le vocabulaire
% moderne sous lequel chercher ce que les feuillets construisent. Le manifeste
% du site (`npm run manifest`) les extrait pour étiqueter la cote entière —
% cette ligne est la source unique des tags, il n'y a pas de fichier à côté.
\newcommand{\keywords}[1]{\par\smallskip\noindent
  {\footnotesize\textsc{Keywords} --- \textit{#1}}\par}

% --- L'appareil critique ----------------------------------------------------

% Le numéro de feuillet, en marge. C'est l'ancre qui permet de régler un
% désaccord sur une formule en regardant *un* feuillet plutôt qu'une chemise
% de six cents. Le site s'en sert aussi pour tourner les pages du fac-similé
% quand on fait défiler la transcription.
\newcommand{\leaf}[1]{%
  \marginnote{\footnotesize\color{gray!70}\textsf{#1}}%
  \label{leaf:#1}\ignorespaces}

% Illisible : on ne devine pas. Un mot inventé qui se lit comme les autres est
% le pire résultat possible.
\newcommand{\ill}{\textcolor{red!60!black}{[\,\ldots\,]}}

% Lecture douteuse : le mot est donné, et signalé comme incertain.
\newcommand{\uncertain}[1]{\uline{#1}}

% Ajout de l'éditeur — une lettre manquante, un mot sous-entendu.
\newcommand{\add}[1]{\textcolor{blue!50!black}{[#1]}}

% Biffé par Grothendieck lui-même. Ce qu'il a rayé fait partie du document :
% une rature dit souvent où la pensée a changé de direction.
\newcommand{\struck}[1]{\sout{#1}}

% Note du transcripteur, sur l'état matériel du feuillet.
\newcommand{\note}[1]{\par\smallskip{\small\color{gray!60!black}\textit{[#1]}}\par\smallskip}

% Note marginale de Grothendieck, distincte du corps du texte.
\newcommand{\marginal}[1]{\par\smallskip\noindent\hspace{1em}%
  {\small\color{gray!60!black}#1}\par\smallskip}

% --- Titre ------------------------------------------------------------------

\AtBeginDocument{%
  \begin{center}
    {\large\bfseries Fonds Alexandre Grothendieck --- cote n\textsuperscript{o}~\grothFolder}\\[2pt]
    {\itshape\grothFoldertitle}\\[4pt]
    {\small feuillets \grothFirst--\grothLast\ (lot \grothBatch)}\\[2pt]
    {\footnotesize\color{gray!60!black}Transcription établie d'après le fac-similé numérisé
     par l'Université de Montpellier.\\
     Les lectures incertaines sont soulignées, les illisibles notées [\,\ldots\,],
     les ajouts entre crochets.}
  \end{center}
  \vspace{1em}}

\endinput


\folder{31}
\batch{5}
\pages{81}{100}
\dating{[1963-1973]}
\watermark{Édition de démonstration}
\foldertitle{SGA A [SGA 4] (Résidus de rédaction) : notes manuscrites (s.d.)}

\begin{document}

\section*{Fin du calcul de $\operatorname{cd}_{\ell}$ (pages 81 et 82)}

\page{81}
\note{La page achève le développement des pages 73 à 80 (batch 4).}

\textbf{Question}\quad $X$ schéma noethérien affine,
\[
n = \operatorname*{Sup}_{x \in X}\bigl[\operatorname{cd}_{\ell}(k(x)) +
\dim \mathcal{O}_{X,x}\bigr].
\]
\marginal{$\ell$ premier \ill{} car.\ rés.}
Alors $n$ est aussi le Sup des \ill{} coh.\ des \ill{} affines de $X$. La
question, \ill{} \ill{}
\[
\operatorname{cd}_{\ell}(X) \leqslant n .
\]

Mais il est faux de dire que l'\ill{} \ill{} en \ill{} d'un exemple de \ill{}
\ill{}, $n = 2$ s'il est \struck{\ill{}} \ill{}, alors que
$\operatorname{cd}(X) = 3$ !

\page{82}

\emph{Dim.\ coh.\ d'un schéma affine sur $\mathbb{Z}$}

$X$, $f : X \to S = \operatorname{Spec}(\mathbb{Z})$, avec $\ell \cdot 1_{X}$
inversible i.e.\ $\ell \notin f(X)$, $\ell \ne 2$. \add{Supposons d'abord $X$
\ill{} de dim.\ relative $d$ sur $S$.}
\[
H^{*}(X,F) \Longleftarrow H^{p}\bigl(S_{\ell}, R^{q} f_{*}(F)\bigr)
\]
\note{La lettre du foncteur image directe est tracée ici d'un jambage plus
ample qu'aux lignes suivantes, où elle est nettement $f$.}
\[
R^{q} f_{*}(F)_{\mathfrak{g}} \ne 0 \Longrightarrow q \leqslant
\struck{1 + } d, \qquad
R^{q} f_{*}(F)_{s} \ne 0 \Longrightarrow q \leqslant 1 + d
\]
\note{$\mathfrak{g}$ est le point générique de $S$, $s$ un point fermé ; la
mention marginale à droite de la première ligne est illisible.}

Donc si $q > 1+d$, $E_{2}^{pq} = 0$ ;

si $q = 1+d$, alors $R^{q} f_{*}(F)_{s}$ est concentré \ill{} sur les points
fermés, donc $E_{2}^{pq} \ne 0 \Rightarrow p \leqslant 1$ ;

si $q \leqslant d$ \struck{$+d$}, donc $E_{2}^{pq} \ne 0 \Rightarrow
p \leqslant 2$.

\note{Un croquis du plan $(p,q)$ accompagne ces trois cas, portant les
ordonnées $1+d$ et $d$ ; il n'ajoute rien aux inégalités.}

On trouve donc
\[
\operatorname{cd}_{\ell}(X) \leqslant d + 2 = \dim X + 1 =
\uncertain{\operatorname{cd}'_{\ell} \mathcal{R}(X)}
\]
\ill{} \ill{} si $X$ est concentré \struck{\ill{}} au dessus d'un point
(\uncertain{Hp}) de $S = \operatorname{Spec}(\mathbb{Z})$. Donc on trouve
\ill{} si $X$ affine de type fini, $\ell \cdot 1_{X}$ inversible, et \ill{}
$X_{\mathbb{Q}}$ \struck{devrait} est au dessus d'un corps \struck{tot.}
\add{tot.} premier \ill{}
\[
\boxed{\operatorname{cd}_{\ell}(X) \leqslant \dim X + 1}
\]

\section*{Théorèmes 3.1 et 3.2 du typescrit (page 83)}
\note{Feuillet du même typescrit que la page 71 (batch 4), corrigé de sa main :
les lettres de groupe et plusieurs indices, que la machine n'avait pas, sont
encrés à la main ; on les transcrit $G$, $H$ comme au batch 4. Les ratures
courtes sont de la machine (frappe sur frappe), les insertions et les
suppressions signalées ci-dessous sont de sa main.}

\page{83}

\textbf{Théorème 3.1.} Soient $S$ un préschéma affine, $G$ un $S$-groupe de
type multiplicatif, $F$ un faisceau quasi-cohérent sur $S$, sur lequel $G$
opère (Exp 1, 4.7.). Alors \struck{pour tout $i > 0$, la cohomologie de
Hochschild} on a
\[
H^{i}(G,F) = 0 \qquad \text{pour } i > 0,
\]
où $H^{i}$ est la cohomologie de Hochschild étudiée dans Exp.\ 1, 5.

En effet, d'après loc.\ cit., 5.3., la cohomologie de Hochschild s'explicite,
en termes des anneaux affines $A$ de $S$ et $B$ de $G$, et \add{du} module $M$
sur $A$ définissant $F$, comme cohomologie d'un complexe \struck{$C^{*}(G,F)$
dont la} de $A$-modules $C^{\bullet}(G,F)$, dont la formation commute \add{au}
changement de base $A \to A'$. Par suite, pour un changement de base
$A \to A'$ avec $A'$ plat sur $A$, on aura
\[
H^{i}(G',F') = H^{i}(G,F) \otimes_{A} A',
\]
et par suite, si on suppose même $A'$ fidèlement plat sur $A$, pour prouver que
$H^{i}(G,F)$ est nul, il suffit de prouver qu'il en est ainsi de
$H^{i}(G',F')$. Cette remarque nous ramène à vérifier 3.1.\ dans le cas où $G$
est diagonalisable, dans ce cas cela a été prouvé dans Exp 1, 5.3.3.

Utilisant les résultats de Exp.\ III, nous allons tirer de 3.1.\ diverses
conséquences de nature géométrique :

\textbf{Théorème 3.2.} Soient $S$ un préschéma affine, $S_{0}$ un sous-schéma
affine défini par un Idéal $J$ de carré nul, $G$ un groupe de type
multiplicatif sur $S$, $H$ un préschéma en groupes quelconque sur $S$,
$u, v : H \to G$ deux homomorphismes de groupes tels que $u_{0} = v_{0}$, où
$u_{0}, v_{0} : H_{0} \to G_{0}$ sont les homomorphismes déduits de $u, v$ par
changement de base $S_{0} \to S$. Alors il existe un \add{$g \in G(S)$} tel que
$v = \operatorname{int}(g) \cdot u$ \add{et $g_{0} = e$}.

Cela résulte de Exp III et de 3.1.

\section*{Conjectures sur la cohomologie des schémas affines (page 84)}

\page{84}

\emph{Conjectures sur la cohomologie des schémas affines}

$X$ un (« bon ») schéma affine, intègre, \ill{} gén.\ $a$, \ill{} $\ell$ un
\ill{} premier.

$1^{\circ})$ Si $\ell \cdot 1_{X}$ inversible, alors
\[
\operatorname{cd}_{\ell}(X) \leqslant \operatorname{cd}_{\ell}(k(x))
\]
\marginal{[et est \ill{} une égalité \ill{} \ill{} \ill{} plus ?]}

$2^{\circ})$ Si $\ell \cdot 1_{X}$ non inversible et non nilpotent, \ill{} a
\[
\operatorname{cd}_{\ell}(X) \leqslant
\operatorname{Sup}\bigl(\operatorname{cd}_{\ell}(k(x)),\;
\nu(V(\ell)) + 3\bigr)
\]
où on pose, pour \ill{} \ill{} un schéma $Z$ de car.\ $\ell$,
\[
\nu(Z) = \operatorname*{Sup}_{z \text{ maximaux dans } Z} \nu(k(z))
\]
et pour un corps $k$ de car.\ $\ell$, $\nu(k)$ est \uncertain{défini par}
$\ell^{\nu(k)} = [k : k^{\ell}]$.

$1\;\mathrm{bis})$ Sous les conditions de $2^{\circ}$, lorsque $\exists\;
\ell'$ tel \struck{que $\ell' \cdot 1_{X}$ inversible i.e.\ $X$ sur
$\operatorname{Spec}(\mathbb{Z}_{\ell'})$ ($S = \operatorname{Spec}
\mathbb{Z}$)} \struck{alors \ill{} \ill{} quelconque \ill{} dans le \ill{}}
\add{$A$} $(\nu(V(\ell)) + 3)$ \ill{} $\nu(V(\ell)) + 1$ \ill{}
\struck{$\nu(V(\ell))$} de $2^{\circ}$, lorsque \struck{Supposons $\exists\;
\ell'$} prenant \ill{} que $\mathbb{Z}_{\ell'}$ inversible. Alors
\[
\operatorname{cd}_{\ell}(X) \leqslant \operatorname{cd}_{\ell}(k(x))
\]
\note{Le passage $1\;\mathrm{bis})$ est repris trois fois sur la page, chaque
reprise raturant la précédente ; la prose de liaison n'est pas récupérable et
une note filant la marge gauche, écrite en travers, l'est encore moins.}

$4^{\circ})$ $A$ bon anneau local hensélien intègre, car.\ résiduelle $\ell$,
$\operatorname{car} K = 0$. Soit $U$ \ill{} ouvert \ill{} de
$X = \operatorname{Spec} A$. Alors
\[
\operatorname{cd}_{\ell}(U) \leqslant \dim A + \nu(k) +
\operatorname{cd}_{\ell}(k)
\]

\section*{Densité schématique : Corollaire 4.6 (page 85)}

\page{85}

Notons en passant le résultat suivant, qui servira dans un exposé ultérieur :

\textbf{Corollaire 4.6.} Soient $X$ un $S$-préschéma plat, $U$ une partie
ouverte de $X$, on suppose que pour tout $s \in S$, $U_{s}$ est
schématiquement dense dans $X_{s}$. Supposons de plus que $X$ est localement
noethérien, ou $X$ localement de présentation finie sur $S$ et $U$
rétrocompact dans $X$ (ou ce qui revient au même, $U$ de présentation finie
sur $X$ ; il suffit par exemple $X$ et $U$ de présentation finie sur $S$).
Alors $U$ est schématiquement dense dans $X$.

Le cas $X$ localement noethérien n'est mis que pour mémoire, il est contenu
dans 4.3. Dans le deuxième cas envisagé, on peut évidemment supposer $S$ et $X$
affine, alors $X$ et $U$ sont de présentation finie sur
$S = \operatorname{Spec}(A)$. \add{$A$ est limite inductive de ses
sous-$\mathbb{Z}$-algèbres $A_{i}$ de type fini.} Le « procédé breveté » déjà
utilisé (explicité en long et en large dans EGA IV 8) montre qu'il existe un
\struck{schéma noethérien affine sous-anneau} $i$, un schéma affine $X_{i}$ sur
$S_{i} = \operatorname{Spec}(A_{i})$ et un ouvert $U_{i}$ dans $X_{i}$, dont
$X$, $U$ se déduisent par changement de base \struck{si} $S \to S_{i}$.
\struck{Pour tout $j \geqslant i$,} soit $E_{i}$ la partie de $S_{i}$ formée
des $s \in S_{i}$ tels que $U_{i,s}$ soit schématiquement dense dans
$X_{i,s}$. Cette condition signifie aussi que $U_{i,s}$ contient l'ensemble
$\operatorname{Ass} \mathcal{O}_{X_{i,s}}$ des points « associés » au faisceau
structural de $X_{i,s}$, et les résultats de EGA IV 9 montrent alors que cette
condition est de nature constructible, i.e.\ que $E_{i}$ est une partie
constructible de $S_{i}$. D'ailleurs, la question envisagée sur les fibres est
invariante par extension du corps de base, ce qui prouve en vertu de
l'hypothèse faite que \struck{que pour tout $j$, \ill{}} l'image inverse de
$E_{i}$ par $S \to S_{i}$ \struck{\ill{}} est $S$, et entraîne par EGA IV 8
qu'il existe un $j \geqslant i$ tel que l'image inverse \struck{de $E_{j}$} de
$E_{i}$ dans $S_{j}$ soit $S_{j}$, i.e.\ que pour tout $s \in S_{j}$,
$U_{j,s}$ est schématiquement dense dans $X_{j,s}$. En vertu de 4.4.\ appliqué
à $(S_{j}, X_{j}, U_{j})$ et au changement de base $S \to S_{j}$, il s'ensuit
que \uncertain{$U$} est schématiquement dense dans $X$. On notera d'ailleurs
qu'on n'utilise ici 4.4.\ que dans le
\note{Le feuillet s'arrête en fin de phrase.}

\section*{Anneaux locaux henséliens à corps résiduel nul (page 86)}

\page{86}

\emph{N.B.} Pour les bons anneaux locaux henséliens de \struck{\ill{}} car.\
résiduelle nulle, tout se simplifie grâce à Artin.

\section*{Corollaires 4.7 et 4.9, Remarque 4.9 (page 87)}

\page{87}

\textbf{Corollaire 4.7.} Soient $X$ un $S$-préschéma plat, $U$ une partie
ouverte de $X$, \ill{} les conditions de 4.6. Alors :

a) pour tout $S$-préschéma plat $X'$, $U \times_{S} X'$ est schématiquement
dense dans $X \times_{S} X'$ ;

b) \ill{} \ill{} \ill{} \ill{} \ill{} \ill{}

\textbf{Corollaire 4.9.} \ill{} \ill{} \ill{} \ill{} \ill{} \ill{} \ill{}
\ill{} \ill{} \ill{}

\textbf{Remarque 4.9.} \ill{} les conditions de 4.6 \ill{} \ill{} \ill{}
\ill{} \ill{} \ill{}
\note{Les deux corollaires portent le même numéro 4.9 ; le chiffre du premier
(4.7) est réencré par-dessus la frappe. On transcrit les numéros tels qu'ils
figurent, sans les régulariser. Le corps de ces trois énoncés est en grande
partie effacé sous les deux longues diagonales d'annulation.}

\section*{Section hyperplane et théorème de Lefschetz (pages 89 et 90)}

\page{89}

\emph{Application}\quad Soit $X$ schéma projectif régulier sur un corps $k$
\add{sép.\ clos} \struck{\ill{}}, $Y$ une section hyperplane régulière, $F$ un
faisceau de $\ell$-\struck{tor}sion loc.\ constant sur $X$ [$\ell$ premier
$\ne$ la car.\ de $k$]. On suppose $\dim X = n$. Considérons $U = X - Y$ et

\begin{tikzcd}[column sep=small, row sep=small, nodes={font=\scriptsize}]
H^{i-1}(U,F) \arrow[r] & H^{i}_{Y}(X,F) \arrow[r] \arrow[d, no head, "\wr"] & H^{i}(X,F) \arrow[r] & H^{i}(U,F) \arrow[r] & H^{i+1}_{Y}(X,F) \arrow[d, no head, "\wr"] \\
 & H^{i-2}(Y, F_{Y} \otimes \check{T}_{\ell}) & & & H^{i-1}(Y, F_{Y} \otimes \check{T}_{\ell})
\end{tikzcd}

et utilisons le fait que
\[
H^{i}(U,F) = 0 \quad \text{si } i > n .
\]

On trouve le

\emph{Th.\ de Lefschetz}
\[
H^{i-2}(Y, F_{Y} \otimes \check{T}_{\ell}) \longrightarrow H^{i}(X,F)
\]
est surjectif si $i \geqslant n+1$, bijectif si $i \geqslant n+2$.

\textbf{Corollaire}\quad Supposons $X$ connexe \add{régulier} (de dim $n$),
alors
\[
H^{2n}(X, \mu_{N}^{\otimes n}) \simeq \mathbb{Z}/N\mathbb{Z}
\]
pour tout \ill{} $N$ premier à \uncertain{la} car.

\page{90}

\emph{Dém.}\quad Récurrence sur $n$, \ill{} \ill{} [$n = 0$ trivial, $n = 1$
c'est connu, \ill{} $n \geqslant 2$], on a \ill{} $2n \geqslant n+2$ d'où on
\ill{} \ill{} le \ill{} de Lefschetz
\[
H^{2n}(X, \mu_{N}^{\otimes n}) \simeq H^{2n-2}(Y, \mu_{N}^{\otimes n} \otimes
\check{T})
\]
\note{Une accolade sous $\mu_{N}^{\otimes n} \otimes \check{T}$ l'identifie à
$\mu_{N}^{\otimes (n-1)}$, suivie de « o.k. ».}

\textbf{Corollaire 2}\quad (Moyennant \ill{})
\[
H^{j}(X,F) \longrightarrow H^{j}(Y,F)
\]
est injectif si $j \leqslant n-1$, bijectif si $j \leqslant n-2$.

Résulte du th.\ de Lefschetz \uncertain{72} par simple
\struck{transposition}, à valeurs dans $\mathbb{Q}/\mathbb{Z}$.

\section*{Théorème de Spécialisation : Corollaire (pages 92 à 98)}
\note{Le titre est de sa main, porté sur la page 91, feuillet par ailleurs
vierge qui ouvre le développement suivant ; ce feuillet ne reçoit donc pas de
numéro de page. Les pages qui suivent sont écrites vite : les énoncés
passent, la prose qui les relie est en grande partie perdue.}

\page{92}

\textbf{Lemme 1}\quad $A$ anneau local noeth.\ \struck{\ill{}} \add{quotient
d'un régulier}, \struck{formellement normal}, $t \in \mathfrak{m}_{A}$
\add{un élément $A$-régulier,} \add{le complété $t$-adique de $A$ un anneau}
\struck{$A/tA$} \struck{formellement} \struck{normal} (a fortiori
\struck{formellement} normal), et $\dim A \geqslant 3$. Soit \struck{$B$ finie
\ill{} $A$}, $B$ « extension » \add{finie} \uncertain{normale} \struck{de}
$A$, et \add{supposons que} $B/tB$ sur $A/tA$ soit \struck{\ill{}} un
revêtement étale \struck{hors de l'origine} \add{hors de l'origine}. Alors $B$
est un revêtement étale \struck{hors de} \ill{}

\emph{N.B.} L'hyp.\ indique que $B$ est étale sur $A$ en les points de
$V(t) - \text{origine}$ (à vérifier : la normalité de \ill{} $A$ et $B$
\ldots) \struck{\ill{} \ill{} \ill{} \ill{} \ill{} existe \ill{}}

\emph{Dém.}\quad Soit $X = \operatorname{Spec} A$, \struck{soit} $X' = X -
\{a\}$, $U$ l'\ill{} \ill{} \ill{} \ill{} que $B$ \add{y} soit étale, alors
\struck{[on a $U \ne X$]} $(X - U) \cap V(t) \subset \{a\}$ \ill{} donc
$\dim(X - U) \leqslant 1$. On a \ill{} un revêtement étale \struck{\ill{}}
$U$, et il faut montrer qu'il se \ill{} en un rev.\ étale de $X'$.
\struck{\ill{}} \ill{} points de $X' - U$, \ill{} $\mathcal{O}_{X,x}$
$\geqslant 2$ [car $\dim \mathcal{O}_{X,x} \geqslant 2$, \ill{} $\geqslant
2$], car \ill{} \ill{}
\marginal{condition des chaînes}

\page{93}
\note{Feuillet très raturé, dont seuls des fragments se laissent lire.}

$(2)$ \ill{} \ill{} $i_{*}(\mathcal{A})$ sur $X'$ \ill{} \ill{}
[$i : U \to X'$] i.e.\ \ill{} canonique, et l'\uncertain{Algèbre} \add{sur $U$}
\ill{} qui \ill{} $V/U$]. Donc, \ill{} \ill{} \ill{}, d'\ill{} \ill{} \ill{}
\ill{} que $X_{1} \to X'$ \ill{} \ill{} $A \to$ \ill{} \ill{}
\struck{\ill{}} duplique $t$-adique.

\struck{De plus \ill{} \ill{} \ill{} complétion des complétés}

\ill{} \ill{} un revêtement \emph{trivial} de $A$. \ill{} \ill{}
\struck{duplique} de $B$ sur $A$. Donc \emph{partout} \ill{} $X'$ \ill{}
$X_{1}$ \ill{} \ill{} \ill{} $B$, \ill{} $X' \times [1,\ill{}]$. Je \ill{}
\marginal{à vérifier soigneusement}
\note{La marge gauche porte en outre un petit carré tracé à la main, ses côtés
étiquetés $A/tA$ ; il schématise le passage au quotient et n'ajoute rien aux
formules.}

\page{94}

\textbf{Lemme 2}\quad Soit $f : X \to Y$ \add{lisse} \struck{simple} un
morphisme \ill{} à fibres de dim 1, $s$ une section de $X$ sur $Y$,
\struck{\ill{}} $Y$ \add{intègre}, \ill{} avec $X$ normal. Soit
\[
X' \longrightarrow X
\]
un rev.\ \ill{} galoisien de groupe $\Gamma$ premier aux car.\ résiduelles de
$Y$, \ill{} \ill{} d'un \ill{} \ill{} \ill{} \ill{} \ill{} $s^{-1}(X')$ \ill{}
\ill{} \ill{}. Alors \ill{}

Supposons de plus que \struck{\ill{} il existe} \ill{} \ill{} \ill{} $B$ \ill{}
$Y$, \ill{} \ill{} \ill{} \ill{} $X - f(Z)$.
\note{Le reste de la page est un palimpseste dont la prose n'est pas
récupérable.}

\page{95}

\ill{} que $g(x)$ \ill{} \ill{} diviseurs de \ill{} \ill{} ramené au cas
\ill{} d'un anneau \ill{} \struck{définie localement par les} \ill{}
certain]. On est ainsi \ill{} \ill{} \ill{} $Y$ \struck{le} spectre \ill{}
d'un anneau de v.\ d.\ discrète, et \ill{} \ill{} [par ex.\ en \ill{}
$\mathcal{O}_{X,g(x)}$].

\textbf{Lemme 3}\quad Soit $A$ un anneau local régulier et de dim 2, $(x,y)$
un syst.\ régulier de paramètres de $A$, $B$ une « extension » \emph{normale}
de $A$, \struck{telle que \ill{}} \ill{}, (i) \ill{} \ill{} galoisien d'ordre
premier \ill{} : \ill{} la car.\ résiduelle de $A$, (ii) $B \to$ \ill{}
l'\ill{} de ramification \ill{} constante dans $V(x)$, (iii) \ill{} \ill{}
complétés en \ill{} \ill{} $y$ de $V(y)$.
\note{Un croquis accompagne l'énoncé : un cercle traversé de deux diamètres qui
se coupent en un point marqué, l'un étiqueté $V(x)$, l'autre $V(y)$ — les deux
diviseurs et leur point d'intersection.}

\page{96}

\ill{} (Abhyankar), \ill{} \ill{} $\Gamma = \mathbb{Z}/n\mathbb{Z}$, $B$ \ill{}
intègre. Alors il \ill{} \ill{} (i) \ill{} (ii) \ill{} \ill{} $A' = A/yA$,
\ill{} \ill{} $B$ \ill{} \ill{}
\[
B \simeq A[T]/(T^{n} - x), \qquad B \otimes_{A} A' = A'[T]/(T^{n} - x'),
\]
\ill{} $x'$ \ill{} l'image de $x$ dans $A'$, \ill{} \ill{} \ill{} \ill{}
\ill{} \ill{} (iii).

\struck{Corollaire}

\textbf{Théorème}\quad Soit \struck{$A$ \ill{} \ill{} noeth.}
\struck{Ceci \ill{} \ill{}} $f : X \to Y$ \struck{$Y$ tor \ill{}} un
\struck{morphisme} \add{morphisme lisse,} \struck{\ill{}} alors $f$ est loc.\
1-asphérique \ill{} pour les groupes finis d'ordre premier aux car.\
résiduelles de $Y$.

\page{97}

\emph{Dém.}\quad On est \ill{} ramené au cas où $Y$ \add{noeth.}
\struck{strictement local}, et par \ill{} \ill{} au cas \ill{} $X \to$ \ill{}
de dim relative 1 sur $Y$. \ill{}

\emph{Applications}\quad $1)$ Soit $f : X \to Y$ un morphisme pr.\ quot.\
\add{lisse} \ill{} \ill{} $R^{i} f_{*}(F)$ localement constant : \ill{}
extension \ill{} \ill{} lisse \ill{} le bon.

\struck{$2)$ Soit $X$}

$2)$ $X$ projectif sur \ill{} \ill{}, \ill{} $k'/k$ ext.\ sép.\ \ill{},
\ill{} \ill{} dernier \ill{} $X$ \ill{} \ill{}
\[
H^{i}(X,F) \;\overset{\sim}{\longrightarrow}\; H^{i}(X',F')
\]
(Étale \ill{} $H^{1}$ et $H^{2}$ \ill{} \ill{} \ill{} \ill{} tifs \ill{}).

$3)$ Si $X$ lisse \ill{} $S$, $Y$ \ill{} \ill{} \ill{} lisse, \ill{} \ill{}

\page{98}

\textbf{Corollaire}\quad Soit \struck{\ill{}} $Y$ \ill{} \ill{} de $X$,
\struck{\ill{} \ill{}} \ill{} équivalentes :
\[
h : V \longrightarrow S, \qquad i \ill{}
\]
\[
F' = \operatorname{Hom}(F, I_{X/S}) \struck{\ill{}}
\]

\begin{itemize}
\item (ii) $R^{i} h_{!}(F'|V)$
\item (i) $R^{i} h_{*}(F|V)$
\item (iii) $R^{i} g_{*}(\struck{F}\;F'|Y)$
\end{itemize}
\note{Une accolade relie (ii) et (i) ; à droite d'elles il écrit « loc.\
constantes pour tt \ill{} ».}

Conditions équivalentes : $Y$ \ill{} \struck{lisse} \ill{}, \ill{} $X$ \ill{}
$S$ \ill{} \ill{} \ill{} variations sur \ill{} $S'$ \ill{} \ill{} $S$, \ill{}
[\ill{} \ill{} ??].
\marginal{Rq.\ \ill{} \ill{} $H^{q}(F)$ \ill{} \ill{} \ill{} \ill{}}

\section*{Note sur la suite spectrale de $R^{i} h_{*}$ (page 100)}

\page{100}
\note{Note de sa main portée sur un feuillet dactylographié sans rapport, que
le scan présente à l'envers ; l'écriture, elle, est à l'endroit.}

\emph{N.B.} \ill{} \ill{} \ill{} \ill{} \ill{} $V$ \ill{} $G_{1}$ \ill{} :
pr.\ l.\ \ill{} \ill{} \ill{} \ill{}
\[
R^{i} h_{*}(G_{1}) \Longleftarrow \ill{}
\]
\ill{} $h : V \to X$ \ill{} \ill{} \ill{}
\[
\ill{} \; R^{i} h_{*}(G).
\]
\ill{} \struck{\ill{}} spectrale ;
\[
c) \Longleftarrow R^{i} g_{*}\bigl(R^{i} i_{*}(G')\bigr), \qquad
R^{i} g_{*}\bigl(R^{i} i_{*}(G)\bigr)
\]
\ill{} \ill{} \ill{} \ill{}
\note{Les deux exposants de la suite spectrale sont écrits $i$ tous les deux ;
on les transcrit tels quels sans les distinguer en $p$ et $q$.}

\end{document}
